% !TEX root = ../main.tex
\chapter{Op Registry and Schema}
\label{ch:op_registry}

Operations are the unit of work in \Lucid: every numerical
computation the framework performs is the application of one or
more operations to one or more tensors. There are roughly 260 of
them in the current release, ranging from arithmetic primitives
(\code{add}, \code{mul}) to neural-network composites
(\code{layer\_norm}, \code{multi\_head\_attention}) to control-flow
constructs (\code{where}, \code{masked\_select}). Each carries
metadata that the dispatcher, the autograd engine, the AMP system,
the compile path, and the introspection tools all consult.

The metadata lives in the op registry, the topic of this chapter.
The registry is a single source of truth: a contributor who adds a
new operation declares its schema in one place and the rest of the
framework picks it up automatically. The registry is also what
makes \Lucid's compile path tractable: the trace IR can be
serialised, the emitter table can be code-generated, and the
backward-pass topology can be derived from the registry's
declarations without inspecting individual op implementations.

This chapter describes the schema, the registry, the binding-code
generation that consumes the registry, and the protocol for adding
a new operation.

\section{Overview}
\label{sec:op_registry:overview}

An \emph{op schema} is a structured record describing one
operation. It contains the op's name, its inputs and outputs, its
backward behaviour, its AMP policy, its determinism guarantees, its
complexity class, and a small set of derived properties (whether it
is differentiable, whether it is a view-producer, whether it has a
constant-shape output).

The \emph{op registry} is the in-process table that holds every
op's schema. It is populated at engine initialisation by
self-registering schema declarations in the per-op C++ files; the
registry is then consulted by:

\begin{itemize}
  \item The dispatcher (\chref{ch:backend_dispatch}), which uses
    the schema to validate argument shapes and dtypes before calling
    the backend.
  \item The autograd engine (\chref{ch:autograd_engine}), which
    uses the schema's backward declaration to construct the
    appropriate \code{grad\_fn} node.
  \item The AMP system (\chref{ch:amp}), which uses the schema's
    AMP policy to decide whether to cast inputs to FP16 or keep
    them in FP32.
  \item The compile path (\chref{ch:op_emitters}), which uses the
    schema to look up the appropriate \MPSGraph\ emitter.
  \item The introspection tools (\code{lucid.\_C.list\_ops()},
    documentation generators, the docs site), which consult the
    registry to enumerate the surface.
\end{itemize}

This many consumers might suggest a complex schema, but in practice
the schema is small: about a dozen fields, all of them simple
types. The discipline is keeping the schema lean enough that adding
a new op is straightforward.

\section{The OpSchema Structure}
\label{sec:op_registry:schema}

The schema is declared in \code{lucid/\_C/ops/registry/OpSchema.h}:

\begin{lstlisting}[style=lucidcpp,language=C++]
struct OpSchema {
  std::string             name;            // canonical op name
  uint32_t                version;         // schema version
  ArgSpec                 inputs[8];       // input tensor specs
  uint8_t                 n_inputs;
  ArgSpec                 outputs[4];      // output tensor specs
  uint8_t                 n_outputs;
  BackwardSpec            backward;        // backward construction
  AmpPolicy               amp;             // FP16/FP32 cast rule
  ComplexityClass         complexity;      // O(1), O(N), O(N log N), O(N^2)
  Determinism             determinism;     // deterministic / not
  uint16_t                flags;           // packed boolean properties
};
\end{lstlisting}

\subsection{Name and version}
\label{ssec:op_registry:name_version}

\code{name} is the canonical op name --- the string that appears
in user-facing API (\code{add}, \code{matmul}, \code{layer\_norm})
and in the trace IR. Names are lowercase, snake-case, and follow
\refframework-compatible conventions where parity matters.

\code{version} is a 32-bit integer that increments when the op's
behaviour changes in a backward-incompatible way. This matters for
checkpoint portability: a checkpoint saved with version 2 of
\code{batch\_norm} cannot be loaded into a process running version
3 without an explicit migration. The version field is the
mechanism the serialisation layer (\chref{ch:serialization}) uses
to refuse incompatible loads.

\subsection{ArgSpec for inputs and outputs}
\label{ssec:op_registry:argspec}

\code{ArgSpec} describes one tensor argument:

\begin{lstlisting}[style=lucidcpp,language=C++]
struct ArgSpec {
  std::string  name;        // parameter name in user API
  DtypeSet     allowed;     // permitted dtypes (bitmask)
  uint8_t      rank_min;    // minimum rank
  uint8_t      rank_max;    // maximum rank (255 = unbounded)
  bool         optional;    // may be absent
  bool         out;         // is this an output (vs input)
};
\end{lstlisting}

The \code{name} matches the Python-side argument name, so the
dispatcher can rebind keyword arguments to positional slots without
maintaining a separate alias table. The \code{allowed} bitmask
enumerates the dtypes the op accepts; the dispatcher raises
\code{LucidError} if an input dtype is not in the set. The
\code{rank\_min} / \code{rank\_max} pair bounds the input rank;
arbitrary ranks use \code{rank\_max = 255}.

\subsection{BackwardSpec}
\label{ssec:op_registry:backward}

\code{BackwardSpec} describes how the op's backward pass is
constructed:

\begin{lstlisting}[style=lucidcpp,language=C++]
enum class BackwardKind {
  None,           // not differentiable
  Inline,         // backward defined in C++
  Custom,         // user-supplied via Function class
  Composite,      // built from forward primitives
};

struct BackwardSpec {
  BackwardKind     kind;
  BackwardFn*      inline_fn;     // for Inline kind
  SavedTensorList  saved;         // which forward tensors to save
};
\end{lstlisting}

Most ops use \code{Inline}: the backward is a known function with
a known shape. The function takes the saved tensors, the output
gradient, and the input shapes, and returns the input gradients.
The autograd engine sets up the \code{grad\_fn} node at forward
time by reading \code{inline\_fn} from the schema.

\subsection{AmpPolicy}
\label{ssec:op_registry:amp}

\code{AmpPolicy} encodes the op's behaviour under mixed precision:

\begin{lstlisting}[style=lucidcpp,language=C++]
enum class AmpPolicy {
  PromoteToWidest,    // default: keep widest input dtype
  CastToHalf,         // force to FP16 under autocast
  PreserveInputs,     // no casting; preserve user dtype exactly
  ForceFloat32,       // always FP32 (reductions, softmax)
};
\end{lstlisting}

Matrix multiplication uses \code{CastToHalf}: under autocast,
\code{matmul} inputs are cast to FP16. Softmax uses
\code{ForceFloat32}: the reduction accumulates in FP32 even when
inputs are FP16, to prevent overflow. Pointwise arithmetic uses
\code{PromoteToWidest}.

\subsection{Complexity, determinism, and flags}
\label{ssec:op_registry:misc}

\code{ComplexityClass} is an enum
(\code{O1}, \code{ON}, \code{ONLogN}, \code{ON2}, \code{ON3}) that
the profiler uses to label per-op costs.

\code{Determinism} is one of \code{Deterministic} or
\code{Nondeterministic}. The \code{set\_deterministic(True)} mode
(\chref{ch:reproducibility}) refuses to execute nondeterministic
ops.

\code{flags} is a 16-bit bitmask of additional boolean properties:
\code{is\_view\_producer}, \code{is\_inplace}, \code{is\_constexpr},
\code{has\_static\_output\_shape}, \code{requires\_contiguous}.

\section{The OpRegistry}
\label{sec:op_registry:registry}

The registry is a singleton hash table from op name to schema,
implemented in
\code{lucid/\_C/ops/registry/OpRegistry.\{h,cpp\}}:

\begin{lstlisting}[style=lucidcpp,language=C++]
class OpRegistry {
public:
  static OpRegistry& instance();
  void           register_op(OpSchema schema);
  const OpSchema& lookup(const std::string& name) const;
  bool            has(const std::string& name) const;
  std::vector<std::string> all_names() const;
private:
  std::unordered_map<std::string, OpSchema> table_;
};
\end{lstlisting}

\subsection{Self-registration}
\label{ssec:op_registry:self_register}

Each op's C++ file contains a self-registering static initialiser:

\begin{lstlisting}[style=lucidcpp,language=C++]
namespace {
  struct AddRegistration {
    AddRegistration() {
      OpSchema s;
      s.name = "add";
      s.version = 1;
      s.n_inputs = 2;
      s.inputs[0] = {"a", DtypeSet::Numeric, 0, 255, false, false};
      s.inputs[1] = {"b", DtypeSet::Numeric, 0, 255, false, false};
      s.n_outputs = 1;
      s.outputs[0] = {"out", DtypeSet::Numeric, 0, 255, false, true};
      s.backward = {BackwardKind::Inline, &add_backward, {}};
      s.amp = AmpPolicy::PromoteToWidest;
      s.complexity = ComplexityClass::ON;
      s.determinism = Determinism::Deterministic;
      OpRegistry::instance().register_op(std::move(s));
    }
  } _add_reg;
}
\end{lstlisting}

The pattern is verbose but explicit: a reader of
\code{Arithmetic.cpp} sees \code{add}'s full schema in one place,
without chasing a macro definition.

\subsection{Initialisation order}
\label{ssec:op_registry:init_order}

\code{OpRegistry::instance()} uses a Meyer's singleton pattern
(\code{static OpRegistry r; return r;}), which is initialised on
first call. Since the first call is the first op's registration,
the table is constructed before any registration runs --- avoiding
the C++ static-initialisation-order fiasco.

\section{Binding-Code Generation}
\label{sec:op_registry:codegen}

The pybind11 binding for each op
(\chref{ch:python_cpp_binding}) is generated from the registry by
a small build-time script, \code{tools/gen\_bindings.py}, run by
CMake before the binding files are compiled. It reads the
registry, then emits one \code{m.def(...)} line per op into
\code{lucid/\_C/bindings/auto\_bind.gen.cpp}.

The pattern handles the cases that are uniform: name, argument
list, default values, return type. Ops with non-uniform
characteristics (custom argument adaptation) are bound by hand and
the generator skips them.

\section{Adding a New Operation}
\label{sec:op_registry:adding}

The protocol:

\begin{enumerate}
  \item Implement the forward in
    \code{lucid/\_C/ops/<category>/<file>.cpp}.
  \item Implement the backward (inline C++ function or composite
    decomposition).
  \item Add the self-registering schema declaration.
  \item Run \code{tools/gen\_bindings.py} (CMake does this
    automatically).
  \item Add a Python-side unit test exercising forward, backward,
    and AMP behaviour.
  \item Add a parity test against the reference framework.
\end{enumerate}

The framework ships a CLI, \code{python -m tools.new\_op}, that
scaffolds the nine-file pattern for the common case
(\code{retro-p6-new-op-cli}).

\section{Schema Versioning and Checkpoint Portability}
\label{sec:op_registry:versioning}

Schema versions exist to keep checkpoints loadable across
behavioural changes. A checkpoint records the schema version of
every op; loading checks every op's version against the registry;
a mismatch triggers either an automatic migration or refusal.

The version is bumped only when behaviour changes affect numerical
output, the saved-tensor set, or argument interpretation. It is
\emph{not} bumped for pure-internal speed-ups or
backward-compatible argument additions.

Two historical migrations exist: \code{batch\_norm v1$\to$v2}
rescaled \code{running\_var} from biased to unbiased; \code{cross\_entropy
v1$\to$v2} changed the default reduction. Migrations live in
\code{serialization/migrations.py} with tests against legacy
checkpoints.

\section{Introspection}
\label{sec:op_registry:introspection}

The registry is queryable from Python:

\begin{lstlisting}[style=lucid,language=LucidPython]
import lucid
names = lucid._C.list_ops()           # ~260 names
schema = lucid._C.op_schema("add")    # dict
lucid._C.has_op("custom_kernel")      # False
\end{lstlisting}

The introspection API is used by the documentation generator, by
the compile path's emitter-table builder, and by the unit test
suite that verifies every op has a parity test.

\section{The Op Lifecycle}
\label{sec:op_registry:lifecycle}

Following a single op from declaration through user invocation
clarifies how the registry interacts with the rest of the engine.
\figref{fig:op_registry:lifecycle} sketches the path; we walk
through it with \code{lucid.add(a, b)} as the example.

\begin{figure}[t]
\centering
\begin{adjustbox}{max width=\textwidth, center}
\begin{tikzpicture}[
  stepbox/.style={draw, rounded corners=2pt, font=\small\sffamily,
               fill=lucid.info.bg, minimum width=8cm,
               minimum height=0.7cm, align=left,
               inner sep=4pt, inner xsep=8pt},
  arrow/.style={->, >=stealth, thick, color=lucid.accent.dark,
                shorten >=2pt, shorten <=2pt},
  brace/.style={decorate, decoration={brace, amplitude=6pt},
                color=lucid.muted, thick},
  lbl/.style={font=\footnotesize\itshape, color=lucid.muted,
              align=left},
]
  \node[stepbox, fill=lucid.note.bg] (decl) at (0, 6)
    {\textbf{Declare:} write OpSchema in C++};
  \node[stepbox, fill=lucid.note.bg] (build) at (0, 5)
    {\textbf{Build:} static initialiser fires at process start};
  \node[stepbox] (call) at (0, 4)
    {\textbf{Call:} user invokes \code{lucid.add(a, b)}};
  \node[stepbox] (lookup) at (0, 3)
    {\textbf{Lookup:} dispatcher reads schema from registry};
  \node[stepbox] (validate) at (0, 2)
    {\textbf{Validate:} check dtypes / shapes / device};
  \node[stepbox] (dispatch) at (0, 1)
    {\textbf{Dispatch:} route to backend per device};
  \node[stepbox] (backward) at (0, 0)
    {\textbf{Record backward:} BackwardSpec builds grad\_fn};
  \node[stepbox, fill=lucid.ok.bg] (result) at (0, -1)
    {\textbf{Return:} TensorImpl with autograd metadata};

  \draw[arrow] (decl) -- (build);
  \draw[arrow] (build) -- (call);
  \draw[arrow] (call) -- (lookup);
  \draw[arrow] (lookup) -- (validate);
  \draw[arrow] (validate) -- (dispatch);
  \draw[arrow] (dispatch) -- (backward);
  \draw[arrow] (backward) -- (result);

  % Right-side brace + labels grouping build-time vs per-call steps
  \draw[brace] (4.3, 6.35) -- (4.3, 4.65);
  \node[lbl, right] at (4.7, 5.5) {build time\\(one-shot)};

  \draw[brace] (4.3, 4.35) -- (4.3, -1.35);
  \node[lbl, right] at (4.7, 1.5)
    {per-call time\\${\sim}1.5\,\mu$s framework\\overhead on M3~Max};
\end{tikzpicture}
\end{adjustbox}
\caption[Op lifecycle.]{Lifecycle of one \Lucid\ operation. The
schema is declared in the op's C++ file and registered at process
start (top two steps; one-shot). Per-call cost is the dispatcher
path: lookup, validation, backend dispatch, backward-graph
construction. Framework overhead is roughly $1.5\,\mu$s on M3~Max;
the backend kernel cost dominates for non-trivial inputs.}
\label{fig:op_registry:lifecycle}
\end{figure}

The dispatcher's validation step deserves expansion. Given the
schema's \code{inputs[].allowed} dtype bitmask, the validator
rejects any input dtype not in the set. Given the
\code{rank\_min}/\code{rank\_max} bounds, it rejects rank
mismatches. The actual shape compatibility check
(for binary broadcast, for reduction axis bounds) is handled by
the per-op forward function rather than by the dispatcher, because
those constraints are not uniformly expressible in the schema.

\section{Implementation: From Schema to Bound Function}
\label{sec:op_registry:impl_path}

To trace the full chain from a schema declaration to a Python
call, follow the example of \code{lucid.gelu}:

\begin{enumerate}
  \item \textbf{Schema declaration} in
    \code{lucid/\_C/ops/ufunc/Activations.cpp}. The schema names
    the op, its single input dtype set (floating types only), its
    backward function pointer, and its AMP policy
    (\code{PreserveInputs}).
  \item \textbf{Self-registration}: the static initialiser at the
    bottom of the file runs at \code{dlopen} time and inserts the
    schema into \code{OpRegistry::instance().table\_}.
  \item \textbf{Binding generation}: at build time,
    \code{tools/gen\_bindings.py} reads the registry, sees
    \code{gelu}, and emits \code{m.def("gelu", \&ops::gelu,
    py::arg("x"))} into \code{auto\_bind.gen.cpp}.
  \item \textbf{Compile and link}: the generated bindings file is
    compiled into the engine extension module.
  \item \textbf{Python wrapper}: in
    \code{lucid/\_ops/composite/activation.py}, a small wrapper
    \code{def gelu(x): return \_wrap(\_C\_engine.gelu(\_unwrap(x)))}
    exposes the op at \code{lucid.nn.functional.gelu} and
    indirectly at \code{lucid.gelu}.
  \item \textbf{User call}: \code{lucid.gelu(x)} dispatches through
    the Python wrapper, the dispatcher, the backend, and back ---
    end to end, ${\sim} 2\,\mu s$ for the framework overhead plus
    the kernel time.
\end{enumerate}

The chain is mechanical once the schema is right. Most of a
contributor's time on a new op is spent on the backward derivation
and the parity tests, not on the registration machinery.

\section{Schema Statistics}
\label{sec:op_registry:stats}

\tabref{tab:op_registry:stats} breaks down the current registry by
op category.

\begin{table}[t]
\centering
\small
\begin{tabular}{lrrl}
\toprule
Category & Count & Differentiable & Example \\
\midrule
Unary (ufunc)    & 64 & 51 & \code{relu}, \code{exp}, \code{erfinv} \\
Binary (bfunc)   & 52 & 34 & \code{add}, \code{atan2}, \code{logaddexp} \\
Generic (gfunc)  & 78 & 62 & \code{reshape}, \code{gather}, \code{cat} \\
NN kernels       & 38 & 38 & \code{conv2d}, \code{layer\_norm}, \code{lstm} \\
Linalg           & 31 & 24 & \code{cholesky}, \code{svd}, \code{solve} \\
FFT              & 22 & 22 & \code{fft}, \code{rfft}, \code{ifftshift} \\
Einops           & 6  & 6  & \code{rearrange}, \code{reduce}, \code{repeat} \\
Complex          & 7  & 7  & \code{real}, \code{imag}, \code{conj} \\
\midrule
Total            & 298 & 244 & \\
\bottomrule
\end{tabular}
\caption[Op registry composition.]{Counts of operations in the
registry by category. Of the 298 registered ops, 244 are
differentiable (have a non-trivial \code{BackwardSpec}). The
remaining 54 are predicates (\code{isnan}, \code{isfinite}),
argument reductions (\code{argmax}, \code{argmin}), bitwise ops,
and rounding ops --- all returning either bool or integer outputs
for which gradient is undefined.}
\label{tab:op_registry:stats}
\end{table}

\section{Summary and Forward References}
\label{sec:op_registry:summary}

The op registry is the single source of truth for \Lucid's
operation surface. Each op declares an \code{OpSchema} containing
its name, version, input/output specs, backward declaration, AMP
policy, complexity class, and determinism guarantee. Schemas
self-register at engine initialisation; the registry is consulted
by the dispatcher, the autograd engine, the AMP system, the compile
path, and the introspection tools.

\chref{ch:kernel_abstraction} treats the kernel layer that the ops
dispatch through. Subsequent chapters treat the op categories one
by one.

\begin{lucidnote}
For the user: the op registry is invisible. The user calls
\code{lucid.add(a, b)}; the framework consults the registry under
the hood to validate arguments, route to a backend, and construct
the autograd node.
\end{lucidnote}
